\chapter{Sensor Hardware and Integration Reference}
\label{app:hardware}

This appendix is a specification-level reference for OpenFlight's actual
and planned building blocks, drawn from vendor datasheets and application
notes (all linked). Notably, OmniPreSense's own sports application note
AN-029 documents the exact OpenFlight golf configuration and cites the
OpenFlight repository by name as its reference
implementation~\cite{an029}.

\section{OmniPreSense OPS243-A (24\,GHz CW Doppler)}

\subsection{Hardware}
Per the product brief~\cite{ops243brief}: 24.00--24.25\,GHz ISM band,
11\,dBm transmit power (FCC ID 2ALLL243A), patch antenna with
\textbf{20$\degs$ azimuth $\times$ 24$\degs$ elevation} $-3$\,dB
beamwidth (footprint $\sim$0.4\,m wide at 1\,m, 1.8\,m at 5\,m); motion
detection 1--100\,m; speed to 348\,mph at 50\,ksps; accuracy spec
0.5\%; USB CDC + 3.3\,V UART (default 19{,}200\,8N1); 5--24\,V supply,
1.7\,W active.

\subsection{API essentials}
From AN-010~\cite{an010}: sample rate \texttt{S=n} (1--1000\,ksps;
10\,ksps default); buffer 1024/512/256/128 via
\texttt{S>}/\texttt{S<}/\texttt{S[}/\texttt{S(}; zero-padding
\texttt{Xn}/\texttt{X=16}/\texttt{X=32} to a 4096-point FFT. Speed
ceiling and resolution scale with sample rate (10\,ksps
$\to$ 31.1\,m/s at 0.061\,m/s; 50\,ksps $\to$ 155.4\,m/s at
0.304\,m/s per 1024-sample buffer). Output modes: \texttt{OJ} JSON,
\texttt{OT} timestamps, \texttt{OM} magnitudes, \texttt{O=n}
multi-object (to 16), \texttt{OF} post-FFT, \texttt{OR} raw I/Q.
Filters: \texttt{R>}/\texttt{R<} speed, \texttt{R$\pm$} direction,
\texttt{M>} magnitude, \texttt{K+} peak averaging; built-in
cosine-error correction \texttt{\^{}/$\pm$n.n} (0--89$\degs$).
\texttt{A!}\ persists settings to flash.

\subsection{Rolling buffer and triggering}
From AN-027~\cite{an027}: \texttt{G1} enters rolling-buffer mode with a
fixed \textbf{4096-sample I/Q buffer organized as 32 segments of 128
samples}; trigger by software (\texttt{S!}) or a 3.3\,V rising edge on
\textbf{J3 pin 3 (HOST\_INT)}; \texttt{S\#n} sets the pre/post-trigger
split (default 8 $\to$ 1024 pre + 3072 post). At OpenFlight's 30\,ksps
the buffer spans 136.5\,ms with a 208.5\,mph ceiling --- AN-027's
``golf ball setting.'' The app note wires a \textbf{SparkFun SEN-14262
Gate output directly to HOST\_INT} (OpenFlight's exact trigger) and
flags the acoustic-latency budget: sound from 2\,m arrives
$\sim$6.6\,ms late, so the pre-trigger split must cover the
impact-to-trigger gap (their example \texttt{S\#18}).

\subsection{The vendor golf recipe (AN-029)}
AN-029~\cite{an029} specifies: \texttt{S=30} (209\,mph ceiling),
\texttt{S(} 128-sample segments, \texttt{X=32} (4096-point FFT,
0.1\,mph resolution, $\sim$200\,Hz report rate), \texttt{US},
\texttt{R>10} to mask waggle, \texttt{M>10}, \texttt{O2} to report
ball + club for smash factor (gated 1.0--1.50, matching
\cref{sec:smash}), \texttt{K+} averaging, sensor 2--3\,m behind the
ball. Golf balls are rated ``high'' reflectivity, detectable 5--10\,m.

\section{RFbeam K-LD7 (24\,GHz FSK, dual-RX angle)}

Per the datasheet~\cite{kld7}: 24.050--24.250\,GHz FSK (two
frequencies, enabling range via phase difference); EIRP 6\,dBm; 1\,TX +
\textbf{2 I/Q RX patches at 6.223\,mm ($\approx\lambda/2$) spacing} ---
the interferometric baseline of \cref{eq:interferometry}; beam
80$\degs$\,H $\times$ 34$\degs$\,V. Per-frame 256-point complex FFT;
\textbf{angle $\pm$90$\degs$ at 1$\degs$ resolution from the
Rx1--Rx2 phase difference}; distance 5\,cm--100\,m (resolution 5\,cm at
the 5\,m range setting); speed 0.1--100\,km/h --- note the
\textbf{62\,mph speed ceiling}, which confines the K-LD7 to angle and
club-speed work; ball speed must come from the OPS243-A. Frame time at
the 100\,km/h setting is 29\,ms ($\sim$34\,Hz).

UART protocol (115200\,8E1 default, to 3\,Mbaud via \texttt{INIT}):
message types \texttt{RADC} (raw ADC: 256\,I + 256\,Q for Rx1@$f_A$,
Rx2@$f_A$, Rx1@$f_B$ --- the payload OpenFlight's interferometry
consumes), \texttt{RFFT} (spectrum + threshold), \texttt{PDAT} (up to
12 raw targets: distance/speed/angle/magnitude), \texttt{TDAT}
(tracked target), \texttt{DDAT} (flags), \texttt{DONE} (frame counter
--- use it to detect dropped frames); requested via the \texttt{GNFD}
bitfield. Configuration: \texttt{RSPI} max speed, \texttt{RRAI} max
range, \texttt{THOF} threshold offset (10--60\,dB), \texttt{RBFR} base
frequency (three channels for multi-module coexistence --- set the two
OpenFlight units to different channels), \texttt{TRFT} tracking filter,
detection-window bounds (\texttt{MIRA}/\texttt{MARA}/\texttt{MIAN}/
\texttt{MAAN}/\texttt{MISP}/\texttt{MASP}). Positive speed = receding.

\section{TI IWR6843 (60--64\,GHz FMCW MIMO, on order)}

Per the datasheet~\cite{iwr6843}: 60--64\,GHz with \textbf{4\,GHz
chirp bandwidth} ($\sim$3.75\,cm native range resolution);
\textbf{3\,TX / 4\,RX = 12 virtual antennas} (TDM-MIMO); on-chip
C674x DSP + radar hardware accelerator (FFT, log-magnitude, CFAR);
complex-baseband ADC to 12.5\,Msps. The LEVM/ISK-class EVMs give
$\sim$120$\degs$ azimuth FoV with $\sim$15$\degs$ azimuth angular
resolution (8 virtual azimuth antennas) and coarse elevation; the AOP
variant trades resolution for a 130$\degs\times$130$\degs$ field. The
out-of-box mmWave SDK demo streams a TLV point cloud
($x,y,z$, Doppler) at $\sim$10--20\,Hz --- \emph{too slow for a 3\,ms
launch window}; using the IWR6843 for launch measurement requires a
custom chirp/frame configuration (short frames, high Doppler span) and
low-level processing, for which TI's people-tracking labs
(TIDEP-01000/01010) are the closest starting points~\cite{iwr6843}.
TI's forums confirm no golf-specific lab exists. The natural OpenFlight
role: a single sensor that measures range, angle, and radial speed
simultaneously (replacing both K-LD7s) once custom chirp work is done.

\section{Raspberry Pi Global Shutter camera (optical module)}

Sony IMX296 sensor: 1456$\times$1088, 3.45\,\si{\micro\meter} pixels,
true global shutter, exposures to $\sim$30\,\si{\micro\second}, max
$\sim$60\,fps streaming~\cite{pigscam}. The key feature for
\cref{app:optical} is the \textbf{XTR external-trigger pad}: pulse low
to expose (exposure = pulse width + 14.26\,\si{\micro\second}); frame
rate follows the pulse train; multiple cameras on one trigger line are
hardware-synchronized; enable with
\texttt{v4l2-ctl -c trigger\_mode=1} (early boards: remove R11 if Q2
is fitted)~\cite{pigscam}. Continuous 60\,fps cannot capture flight ---
the strobed multi-exposure design of \cref{app:optical} is the correct
use of this sensor.

\section{Simulator integration: GSPro Open Connect}
\label{app:gspro}

GSPro's Open Connect v1~\cite{gspro} is the only fully documented open
launch-monitor protocol and should be OpenFlight's native output
(PiTrac ships it; E6/TruGolf and Foresight FSX require partnership
agreements~\cite{pitrac}). Mechanics: JSON over TCP to
\texttt{127.0.0.1:0921}, launch monitor as client. Minimum shot
message: \texttt{DeviceID}, \texttt{ShotNumber}, \texttt{APIversion},
\texttt{ShotDataOptions\{ContainsBallData, ContainsClubData\}}, and
\texttt{BallData} with the five required fields:
\begin{center}\small
\begin{tabular}{@{}llll@{}}
\toprule
Field & Units & Required & Maps to \\
\midrule
\texttt{Speed} & mph & yes & ball speed (\cref{tab:ballparams}) \\
\texttt{VLA} & deg & yes & launch angle \\
\texttt{HLA} & deg & yes & launch direction \\
\texttt{TotalSpin} & rpm & yes$^{*}$ & spin rate \\
\texttt{SpinAxis} & deg & yes$^{*}$ & spin-axis tilt ($-$ = draw) \\
\bottomrule
\end{tabular}

\smallskip
\footnotesize $^{*}$or \texttt{BackSpin}+\texttt{SideSpin}, related by
\cref{eq:spincomponents}.
\end{center}
Optional \texttt{ClubData} carries \texttt{Speed},
\texttt{AngleOfAttack}, \texttt{Path}, \texttt{FaceToTarget},
\texttt{Loft}, \texttt{Lie}, \texttt{SpeedAtImpact},
\texttt{VerticalFaceImpact}, \texttt{HorizontalFaceImpact},
\texttt{ClosureRate} --- precisely the \cref{tab:clubparams} set, so
the provenance tagging of \cref{ch:implications} carries through
unchanged. Status flags (\texttt{LaunchMonitorIsReady},
\texttt{LaunchMonitorBallDetected}, \texttt{IsHeartBeat}) and the
201 player-info response (handedness, selected club --- useful for
per-club priors, \cref{app:priors}) complete the loop.

\section{Regulatory constants for the physics engine}

USGA/R\&A equipment rules anchor the models of
\cref{ch:impact,ch:flight}~\cite{usgarules}: ball mass
$\le45.93$\,g and diameter $\ge42.67$\,mm (the constants in
\cref{eq:eom}); clubhead characteristic time
$\mathrm{CT}\le239$\,\si{\micro\second} (+18 tolerance)
$\approx$ COR 0.830 --- the $e$ in \cref{eq:ballspeed}; head MOI
$\le5900$\,g\,cm$^2$ (+100) about the vertical CG axis --- the upper
bound on $I_h$ in \cref{eq:gear}; volume $\le460$\,cc. The Overall
Distance Standard (317\,yd + 3 at 120\,mph clubhead / 10$\degs$ /
2520\,rpm; from January 2028, 125\,mph / 11$\degs$ / 2200\,rpm) is a
useful end-to-end sanity envelope for the flight model: a conforming
ball simulated at ALC conditions must not materially exceed
320\,yd~\cite{usgarules}.

\section{Detection theory: CFAR selection}

OpenFlight's current CFAR (SNR $>15$ over a 150-bin DC mask) is a
cell-averaging scheme. The radar literature~\cite{cfartutorial}
distinguishes: CA-CFAR (mean of training cells around the cell under
test, guard cells excluded; threshold $\alpha\hat P_n$ with $\alpha$
set by the desired false-alarm probability) --- optimal in homogeneous
noise; and \textbf{OS-CFAR} (order statistic: the $k$-th ranked
training cell) --- robust when two targets sit close together, at
$\sim$0.5--1\,dB detection loss. The club-then-ball geometry of a golf
shot is precisely the two-closely-spaced-targets case, so OS-CFAR is
the better default for the impact window.
